\documentclass[a4paper,11pt]{article}
\pagestyle{empty}
\setlength{\parindent}{0pt}
\setlength{\parskip}{0.25em}
\setlength{\textwidth}{6.5in}
\setlength{\textheight}{9in}
\setlength{\oddsidemargin}{-0.3in}
\setlength{\evensidemargin}{-0.3in}
\setlength{\topmargin}{-0.4in}
\setlength{\headheight}{0pt}
\setlength{\headsep}{0pt}

\begin{document}

\textbf{Date:} 2026-07-30

\vskip 0.2em
\textbf{To:} Editor-in-Chief, IEEE Transactions on Information Forensics and Security

\vskip 0.2em
\textbf{Subject:} Submission of ``\textit{Decision-Grade Masking}\,''

\vskip 0.6em

Dear Editor,

We respectfully submit the enclosed manuscript for consideration in IEEE TIFS. The work addresses a long-standing question at the intersection of formal verification and side-channel security: when is the equipment-free statistical channel's verdict sound for the glitch-extended probing model?

\textbf{Key contributions.}
\begin{itemize}
\itemsep0.1em
\item A sound-by-construction static lint that, by Theorem~1, certifies when the equipment-free statistical channel's PASS entails the glitch-extended verdict.
\item A \emph{four-tier} decision-grade evidence chain: functional-stage alive, detection-stage alive, masked $d{=}2$ verdict, masked $d{=}1$ baseline. Each tier is bit-exact reproducible on the same Docker image.
\item End-to-end reproducibility in $\approx 4$~minutes wall time on a standard 8-thread workstation, with the full Docker image and configuration supplied.
\item A direct comparison with three contemporary automated verifiers (maskVerif, IronMask, VERICA), clarifying the distinction between exact decision procedures and our sound cross-channel certification.
\end{itemize}

\textbf{Artifact availability.}
For double-blind review, the artifact is hosted at the anonymized URL: \texttt{https://anonymous.4open.science/r/probing-resistant-masked-hardware}. This URL provides a frozen snapshot of the source tree, all 33 PROLEAD batch outputs (Tier~3), five positive-control runs, and a reproducibility recipe. The Docker image is \texttt{prolead-d2-experimental:latest} (4.63~GB; full SHA-256 redacted here and revealed at camera-ready). Source code and experiment results are also supplied as supplementary ZIPs.

\textbf{Decision-grade evidence.}
Across 33 batches of $64{,}000$ simulations each ($17{,}157$ deduplicated probing sets, Bonferroni threshold $6.54$), we observe \textbf{0/17,157 LEAK}. The same Docker image, run on an unmasked first-order replica, yields $1{,}224/1{,}890$ LEAK with worst-case $-\log_{10}p = +\infty$, confirming that the toolchain's detection stage is alive. A masked first-order baseline (legacy macOS run, $14{,}400$ simulations) returns 0/14,400 with worst-case $-\log_{10}p = 3.39$, confirming that the absence of verdicts in Tier~3 is genuine and not an artifact of a stalled detection stage.

\textbf{Submission and double-blind policy.}
The manuscript follows TIFS double-blind policy: author identities, affiliations, and direct repository URLs are anonymized throughout. Camera-ready will disclose author names, the canonical GitHub URL, and the full Docker image SHA-256. Suggested reviewers (anonymized, no conflict of interest) and a full conflict-of-interest disclosure are attached.

We thank the editor and reviewers for their time and consideration. Please direct any correspondence to the corresponding author via the ScholarOne portal.

\vskip 0.5em

Sincerely,\\
\textbf{Anonymous Author(s)}\\
Sejong University

\end{document}
